% ZHOU-SOURCE-PAGE: 136 PRINTED: 120

starting state $0$ (the new toss yields $T$) and we need another expected
$\E[f(n+1)]$ tosses to reach $(n+1)H$. So we have

\[
\E[f(n+1)]=\E[f(n)]+\frac{1}{2}\cdot1+\frac{1}{2}\cdot\E[f(n+1)]
\]

\[
\Rightarrow \E[f(n+1)]=2\cdot\E[f(n)]+2=2^{n+2}-2.
\]

\emph{General Martingale approach:} Let's use $HH\cdots H_n$ to explain a
general approach for the expected time to get any coin sequence by exploring the
stopping times of martingales.\footnote[9]{If you prefer more details about the
approach, please refer to ``A Martingale Approach to the Study of Occurrence of
Sequence Patterns in Repeated Experiments'' by Shuo-Yen Robert Li, \textit{The
Annals of Probability}, Vol. 8, No. 6 (Dec., 1980), pp. 1171--1176.} Imagine a
gambler has \$1 to bet on a sequence of $n$ heads ($HH\cdots H_n$) in a fair game
with the following rule: Bets are placed on up to $n$ consecutive games (tosses)
and each time the gambler bets all his money (unless he goes bankrupt). For
example, if $H$ appears at the first game, he will have \$2 and he will put all
the \$2 into the second game. He stops playing either when he loses a game or when
he wins $n$ games in a roll, in which case he collects \$\(2^n\) (with probability
$1/2^n$). Now let's imagine, instead of one gambler, before each toss a new
gambler joins the game and bets on the same sequence of $n$ heads with a bankroll
of \$1 as well. After the $i$-th game, $i$ gamblers have participated in the game
and the total amount of money they have put in the game should be \$i. Since each
game is fair, the expected value of their total bankroll is \$i as well. In other
words, if we denote $x_i$ as the amount of money all the participating gamblers
have after the $i$-th game, then $(x_i-\ii)$ is a martingale.

Now, let's add a stopping rule: the whole game will stop if one of the gamblers
becomes the first to get $n$ heads in a roll. A martingale stopped at a stopping
time is a martingale. So we still have $\E[(x_i-\ii)]=0$. If the sequence stops after
the $i$-th toss ($i\geq n$), the $(\ii-n+1)$-th player is the (first) player who
gets $n$ heads in a roll with payoff $2^n$. So all the $(\ii-n)$ players before him
went bankrupt; the $(\ii-n+2)$-th player gets $(n-1)$ heads in a roll with payoff
$2^{n-1}$; $\ldots$; the $i$-th player gets one head with payoff $2$. So the total
payoff is fixed and $x_i=2^n+2^{n-1}+\cdots+2^1=2^{n+1}-2$.

Hence, $\E[(x_i-\ii)]=2^{n+1}-2-\E[\ii]=0\Rightarrow \E[\ii]=2^{n+1}-2$.

This approach can be applied to any coin sequences---as well as dice sequences or
any sequences with arbitrary number of elements. For example, let's consider the
sequence $HHTTH$. We can again use a stopped martingale process for sequence
$HHTTH$. The gamblers join the game one by one before each toss to bet on the same
sequence $HHTTH$ until one gambler becomes the first to get the sequence $HHTTH$.
If the sequence stops after the $i$-th toss, the $(\ii-5)$th gambler gets the
$HHTTH$ with payoff

% ZHOU-SOURCE-PAGE: 137 PRINTED: 121

$2^6$. All the $(\ii-6)$ players before him went bankrupt; the $(\ii-4)$th player
loses in the second toss ($HT$); the $(\ii-3)$th player and the $(\ii-2)$th player
lose in the first toss ($T$); the $(\ii-1)$th player gets sequence $HH$ with payoff
$2^2$ and the $i$-th player gets $H$ with payoff 2.

Hence, $\E[(x_i-\ii)]=2^6+2^2+2^1-\E[\ii]=0\Rightarrow \E[\ii]=70$.

\end{solution}

\subsection{Dynamic Programming}

Dynamic Programming refers to a collection of general methods developed to solve
sequential, or multi-stage, decision problems.\footnote[10]{This section barely
scratches the surface of dynamic programming. For up-to-date dynamic programming
topics, I'd recommend the book \textit{Dynamic Programming and Optimal Control} by
Professor Dimitri P. Bertsekas.} It is an extremely versatile tool with
applications in fields such as finance, supply chain management, and airline
scheduling. Although theoretically simple, mastering dynamic programming
algorithms requires extensive mathematical prerequisites and rigorous logic. As a
result, it is often perceived to be one of the most difficult graduate level
courses.

Fortunately, the dynamic programming problems you are likely to encounter in
interviews---although you often may not recognize them as such---are rudimentary
problems. So in this section we will focus on the basic logic used in dynamic
programming and apply it to several interview problems. Hopefully the solutions to
these examples will convey the gist and the power of dynamic programming.

A discrete-time dynamic programming model includes two inherent components:

\noindent\emph{1. The underlying discrete-time dynamic system}

A dynamic programming problem can always be divided into stages with a decision
required at each stage. Each stage has a number of states associated with it. The
decision at one stage transforms the current state into a state in the next stage
(at some stages and states, the decision may be trivial if there is only one
choice).

Assume that the problem has $N+1$ stages (time periods). Following the convention,
we label these stages as $0,1,\ldots,N-1,N$. At any stage $k$, $0\leq k\leq N-1$,
the state transition can be expressed as

\[
x_{k+1}=f(x_k,u_k,w_k),
\]

where $x_k$ is the state of system at stage $k$,\footnote[11]{In general, $x_k$
can incorporate all past relevant information. In our discussion, we only consider
the present information by assuming Markov property.} $u_k$ is the decision
selected at stage $k$; $w_k$ is a random parameter (also called disturbance).

% ZHOU-SOURCE-PAGE: 138 PRINTED: 122

Basically the state of next stage $x_{k+1}$ is determined as a function of the
current state $x_k$, current decision $u_k$ (the choice we make at stage $k$ from
the available options) and the random variable $w_k$ (the probability distribution
of $w_k$ often depends on $x_k$ and $u_k$).

\noindent\emph{2. A cost (or profit) function that is additive over time.}

Except for the last stage $(N)$, which has a cost/profit $g_N(x_N)$ depending only
on $x_N$, the costs at all other stages $g_k(x_k,u_k,w_k)$ can depend on $x_k$,
$u_k$, and $w_k$. So the total cost/profit is

\[
g_N(x_N)+\sum_{k=0}^{N-1}g_k(x_k,u_k,w_k).
\]

The goal of optimization is to select strategies/policies for the decision
sequences $\pi^*=\{u_0^*,\ldots,u_{N-1}^*\}$ that minimize expected cost (or
maximize expected profit):

\[
\begin{aligned}
J_{\pi^*}(x_0)
&=\min_{\pi}\E\left\{g_N(x_N)+\sum_{k=0}^{N-1}g_k(x_k,u_k,w_k)\right\}.
\end{aligned}
\]

\emph{Dynamic programming (DP) algorithm}

The dynamic programming algorithm relies on an idea called the \emph{Principle of
Optimality}: If $\pi^*=\{u_0^*,\ldots,u_{N-1}^*\}$ is the optimal policy for the
original dynamic programming problem, then the tail policy
$\pi_i^*=\{u_i^*,\ldots,u_{N-1}^*\}$ must be optimal for the tail subproblem

\[
\begin{aligned}
\E\left\{g_N(x_N)+\sum_{k=\ii}^{N-1}g_k(x_k,u_k,w_k)\right\}.
\end{aligned}
\]

\emph{DP algorithm:} To solve the basic problem
\[
\begin{aligned}
J_{\pi^*}(x_0)
&=\min_{\pi}\E\left\{g_N(x_N)+\sum_{k=0}^{N-1}g_k(x_k,u_k,w_k)\right\},
\end{aligned}
\]
start with $J_N(x_N)=g_N(x_N)$, and go backwards minimizing cost-to-go function
$J_k(x_k)$:

\[
\begin{aligned}
J_k(x_k)
&=\min_{u_k\in U_k}\E\left\{g_k(x_k,u_k,w_k)
 +J_{k+1}\bigl(f(x_k,u_k,w_k)\bigr)\right\},\\
&\hspace{8em} k=0,\ldots,N-1.
\end{aligned}
\]

Then the $J_0(x_0)$ generated from this algorithm is the expected optimal cost.

Although the algorithm looks complicated, the intuition is straightforward. For
dynamic programming problems, we should start with optimal policy for every
possible state of the final stage (which has the highest amount of information and
least amount of uncertainty) first and then work backward towards earlier stages
by applying the tail policies and cost-to-go functions until you reach the initial
stage.

Now let's use several examples to show how the DP algorithm is applied.

% ZHOU-SOURCE-PAGE: 139 PRINTED: 123

\begin{problem}{Dice game}

You can roll a 6-side dice up to 3 times. After the first or the second roll, if you
get a number $x$, you can decide either to get $x$ dollars or to choose to continue
rolling. But once you decide to continue, you forgo the number you just rolled. If
you get to the third roll, you'll just get $x$ dollars if the third number is $x$
and the game stops. What is the game worth and what is your strategy?
\end{problem}

\begin{solution}
This is a simple dynamic programming strategy game. As all dynamic programming
questions, the key is to start with the final stage and work backwards. For this
question, it is the stage where you have forgone the first two rolls. It becomes a
simple dice game with one roll. Face values 1, 2, 3, 4, 5, and 6 each have a $1/6$
probability and your expected payoff is \$3.5.

Now let's go back one step. Imagine that you are at the point after the second roll,
for which you can choose either to have a third roll with an expected payoff of
\$3.5 or to keep the current face value. Surely you will keep the face value if it
is larger than 3.5; in other words, when you get 4, 5 or 6, you stop rolling. When
you get 1, 2 or 3, you keep rolling. So your expected payoff before the second roll
is $3/6\cdot3.5+1/6\cdot(4+5+6)=\$4.25$.

Now let's go back one step further. Imagine that you are at the point after the first
roll, for which you can choose either to have a second roll with expected payoff
\$4.25 (when face value is 1, 2, 3 or 4) or keep the current face value. Surely you
will keep the face value if it is larger than 4.25; In other words, when you get 5
or 6, you stop rolling. So your expected payoff before the first roll is
$4/6\cdot4.25+1/6\cdot(5+6)=\$14/3$.

This backward approach---called tail policy in dynamic programming---gives us the
strategy and also the expected value of the game at the initial stage, \$14/3.

\end{solution}

\begin{problem}{World series}

The Boston Red Sox and the Colorado Rockies are playing in the World Series finals.
In case you are not familiar with the World Series, there are a maximum of 7 games
and the first team that wins 4 games claims the championship. You have \$100 dollars
to place a double-or-nothing bet on the Red Sox.

Unfortunately, you can only bet on each individual game, not the series as a whole.
How much should you bet on each game so that if the Red Sox wins the whole series,
you win exactly \$100, and if Red Sox loses, you lose exactly \$100?
\end{problem}

\begin{solution}
Let $(i,j)$ represents the state that the Red Sox has won $i$ games and the Rockies
has won $j$ games, and let $f(i,j)$ be our net payoff, which can be negative when we
lose money, at state $(i,j)$. From the rules of the game, we know that there may be
between 4 and 7 games in total. We need to decide on a strategy so that whenever the

% ZHOU-SOURCE-PAGE: 140 PRINTED: 124

series is over, our final net payoff is either +100---when Red Sox wins the\linebreak
championship---or -100---when Red Sox loses. In other words, the state space of the
final stage includes $\{(4,0),(4,1),(4,2),(4,3)\}$ with payoff $f(i,j)=100$ and
$\{(0,4),(1,4),(2,4),(3,4)\}$ with payoff $f(i,j)=-100$. As all dynamic
programming questions, the key is to start with the final stage and work
backwards---even though in this case the number of stages is not fixed. For each
state $(i,j)$, if we bet \$y on the Red Sox for the next game, we will have
$(f(i,j)+y)$ if the Red Sox wins and the state goes to $(\ii+1,j)$, or
$(f(i,j)-y)$ if the Red Sox loses and the state goes to $(i,j+1)$. So clearly we
have

\[
\begin{aligned}
f(\ii+1,j)&=f(i,j)+y,\\
f(i,j+1)&=f(i,j)-y
\end{aligned}
\qquad\Rightarrow\qquad
\begin{cases}
f(i,j)=\dfrac{f(\ii+1,j)+f(i,j+1)}{2},\\[6pt]
y=\dfrac{f(\ii+1,j)-f(i,j+1)}{2}.
\end{cases}
\]

For example, we have
$f(3,3)=\dfrac{f(4,3)+f(3,4)}{2}=\dfrac{100-100}{2}=0$. Let's set up a table
with the columns representing $i$ and the rows representing $j$. Now we have all the
information to fill in $f(4,0)$, $f(4,1)$, $f(4,3)$, $f(4,2)$, $f(0,4)$, $f(1,4)$,
$f(2,4)$, as well as $f(3,3)$. Similarly we can also fill in all $f(i,j)$ for the
states where $i=3$ or $j=3$ as shown in Figure 5.7. Going further backward, we can
fill in the net payoffs at every possible state. Using equation
$y=(f(\ii+1,j)-f(i,j+1))/2$, we can also calculate the bet we need to place at each
state, which is essentially our strategy.

If you are not accustomed to the table format, Figure 5.8 redraws it as a binomial
tree, a format you should be familiar with. If you consider that the boundary
conditions are $f(4,0)$, $f(4,1)$, $f(4,3)$, $f(4,2)$, $f(0,4)$, $f(1,4)$, $f(2,4)$,
and $f(3,4)$, the underlying asset either increases by 1 or decrease by 1 after
each step, and there is no interest, then the problem becomes a simple binomial tree
problem and the bet we place each time is the delta in dynamic hedging. In fact,
both European options and American options can be solved numerically using dynamic
programming approaches.
\end{solution}
